\section{Conclusiones y recomendaciones}

\begin{enumerate}
  \item GymHub responde de forma coherente a la hipótesis de centralizar procesos
  administrativos y deportivos, y el prototipo demuestra que la solución puede
  construirse con las tecnologías seleccionadas.
  \item La existencia de un MVP no demuestra demanda. La viabilidad comercial
  permanece condicionada a la encuesta, la disposición de pago y un piloto con
  un gimnasio local.
  \item El modelo de suscripción es consistente con costos recurrentes de nube,
  soporte y mantenimiento. Un escenario conservador requiere aproximadamente
  tres clientes básicos o dos clientes con ingreso promedio intermedio para
  cubrir únicamente costos mensuales de caja.
  \item La valoración económica debe incorporar horas de trabajo, dominio,
  impuestos, soporte y contingencias antes de concluir que el negocio es
  rentable.
  \item El alcance inicial debe priorizar miembros, membresías, pagos, asistencia,
  planes y progreso. Las funciones avanzadas deben agregarse solo si la evidencia
  muestra valor suficiente.
  \item La privacidad constituye un requisito central porque se procesan datos
  personales, financieros operativos y de actividad física. Las demostraciones
  deben permanecer anonimizadas.
  \item El módulo nutricional debe limitarse a orientación general y no debe
  presentarse como intervención médica o sustitución de profesionales.
  \item La arquitectura requiere trabajo adicional para producción: separación
  por gimnasio, workers persistentes, almacenamiento de media, respaldos y QA
  completo.
  \item La propuesta puede crecer si mantiene una adopción simple, soporte cercano
  y decisiones basadas en evidencia, en lugar de acumular funciones sin uso
  comprobado.
\end{enumerate}

Se recomienda aplicar primero el instrumento de validación, actualizar el análisis
con sus resultados y seleccionar un aliado para un piloto de alcance limitado. En
paralelo, debe cerrarse la brecha multiempresa y documentarse un procedimiento de
seguridad, respaldo y atención. Solo después de medir uso, costo de soporte y
disposición de pago conviene fijar precios comerciales.

\section{Declaración de uso de inteligencia artificial}

Para el desarrollo de GymHub se utilizó inteligencia artificial como herramienta
de apoyo en programación, resolución de dudas técnicas, organización del
documento, revisión de redacción y estructuración de escenarios. El equipo
estudiantil es responsable de revisar, adaptar y validar el código, las fuentes,
los cálculos y el contenido final. La inteligencia artificial no sustituyó la
investigación de campo; por esa razón, los resultados no recolectados permanecen
identificados como pendientes.
